\section{Features and training/serving parity}
\label{sec:features}

\subsection{The feature catalogue}

Every transaction yields a vector of \nFeatures{} features in \nFeatureGroups{} groups: velocity
(counts and sums over trailing windows, and the ratio of a short window's count to a long window's
mean), amount behaviour, temporal (local hour and day from the country pack's UTC offset, night,
month-end salary windows, time since the last transaction), geographic (distance and implied speed
from the last transaction, distance from a home centroid, first transfer to a
counterparty country, and a
lagged, smoothed labelled-fraud rate per map cell), counterparty (novelty, age, distinct senders as
a mule signal, lagged labelled fraud), device and channel, account profile (including days since a
SIM swap), agent-specific, corridor, and a label-free synthetic-identity composite. Appendix~\ref{app:features}
gives each feature's definition, with its windows and constants, as the feature registry states it
at \texttt{m4-complete}; the appendix is generated from the registry and checked against it, and is
not typed by hand. Every label-derived feature uses only labels whose availability time precedes the
transaction.

The composite is the unweighted mean of \nSynthTerms{} terms: low KYC tier, a new account, device
sharing and a volume ramp. This benchmark has no KYC tier and no account opening date
(Section~\ref{sec:limitations}), so the KYC and new-account terms always contribute nothing and only the
device-sharing and volume-ramp terms can contribute; whether the device-sharing term varies on the
evaluated draw has not been measured. On this benchmark \nNoSourceFeatures{} of the features have no
source data and one is constant, so the models train on \nTrainableFeatures{}.

\subsection{Batch and online paths, and a test that can fail}

A model is trained on vectors computed in batch over history and served vectors computed online,
one transaction at a time. If they differ, every offline metric describes a system that does
not exist in production, and the gap is invisible in both. The parity design
(\texttt{docs/ml/training\_serving\_parity.md}) settles the following:
\begin{enumerate}
  \item \textbf{What counts as identical.} The specification asks for agreement to an absolute
        $\nSpecTolerance{}$. As an absolute tolerance that is vacuous on a flag (\nSlackUlps{}
        units in the last place of slack at a magnitude of one) and unsatisfiable on the largest
        week-long sums, where it falls below the spacing of representable float64 values.
        Integral and categorical features must match exactly; real-valued features must satisfy
        $|b - o| \le \nParityRel{} + \nParityRel{}\,|b|$; NaN positions must match exactly
        (ADR 0025). The relative rule is stricter than the specification's for values of magnitude
        below \nParityCrossover{} and looser above it, which is where the absolute rule cannot be
        met.
  \item \textbf{What the paths are fed.} Prefix replay: the online path receives an account's
        transactions one at a time, and the batch path is computed on the prefix ending at each
        one. A batch window that reaches forward in time fails immediately.
  \item \textbf{Which cases are covered}: an empty history, both sides of every window boundary, all
        \nChannels{} channels (so the NaN contract is exercised), label-derived features against
        their availability time, per-account replay, a cold cache flushed mid-replay, and a history
        that spans a configuration change, asserted to contain one before agreement is asserted,
        since the paths agree on every transaction newer than the last change (ADR 0026).
  \item \textbf{Independence, enforced in part.} A parity test is blind to a bug both paths share,
        so a test fails if either path directly imports the other's module. It reads each module's
        own imports, not the whole import graph, and the paths are permitted to share the
        declarative registry, the data types and a set of computing primitives, so a bug in a
        shared primitive is outside what it can detect.
  \item \textbf{Proof the test can fail.} Deliberately wrong online paths (mutations) must be
        caught; a mutation that passes specifies a missing case.
  \item \textbf{The database fallback is a further path.} A database aggregate that serves velocity
        features when the cache is down is another computation of the same quantities and is held
        to the same contract.
\end{enumerate}
The limit is stated rather than hidden: the parity suite only shows that the paths agree on the
inputs they are given. The velocity ratio did not filter its history by account on the batch path;
the online path could not make that mistake, because its state is keyed by account, but every
parity and unit fixture was a single-account history, so nothing exposed the defect. An
index-equivalence test that fed the batch path a mixed corpus found it. Feature computation is
deterministic, byte-identical across batch sizes of \nBatchSizes{}. The online path's latency
(the M3 target is a \nPctlFeature{}th-percentile latency below \nFeatureLatencyTarget{}~ms including the cache fetch) is
\deferredmachine{}; latency percentiles count for the gate only from such a machine (ADR 0010).

\subsection{The single-feature floor}
\label{sec:floor}

The single-column separation control of Section~\ref{sec:dataset} holds for the dataset's columns
and not for the engineered features. Measured on the evaluation draw at \nReScoredRows{} scored rows (the most recent rows of a
\nReCorpusRows{}-row corpus, so each has the corpus behind it as history)
holding \nReScoredFraud{} observed fraud, with categorical encodings folded by whole account,
\nReAboveCeiling{} of the \nFeatures{} features exceed the \nCeiling{} ceiling alone at their point
estimates (Table~\ref{tab:floor}); with intervals as wide as these, \nReAboveCeilingCI{} exceeds it
beyond its interval. The strongest, \texttt{velocity\_ratio\_1h\_vs\_30d}, reaches \nReVelocity{}
$\pm$ \nReVelocityci{}. On the evaluation draw's whole test period the same feature reaches
\nGateFloor{} (Section~\ref{sec:evaluation}); each gate ranking metric's margin is measured against
this feature on that metric's own rows.

The measurement was first taken on an earlier draw, of \nFloorDrawRows{} rows at the same seed,
whose evidence file records no fingerprint; there \nFeaturesAboveCeiling{} features exceeded the
ceiling at their point estimates, the velocity ratio at \nFloorM{}. Re-run on the evaluation draw with
the same command, \nReAboveCeiling{} of them exceed it, and the others fall just below it, with
intervals that still reach it. How many features cross a fixed ceiling therefore depends on the draw
and the sample; in both measurements the velocity ratio is beyond it by more than its interval.

\begin{table}[htbp]
\centering
\caption{The features above the separation ceiling (\nCeiling{}) on the earlier draw,
$\max(\mathrm{AUC}, 1-\mathrm{AUC})$, measured again on the evaluation draw \texttt{d8083dbc}. First
column: the recorded command re-run (\nReScoredRows{} scored rows holding \nReScoredFraud{} fraud,
Hanley--McNeil intervals; \texttt{evidence/\allowbreak round2\_\allowbreak single\_\allowbreak feature\_\allowbreak auc\_\allowbreak d8083dbc.txt}).
Second column: the same feature on the gate's \nTestRows{} test rows with an account-clustered
interval, where the feature is defined on every row (\texttt{evidence/\allowbreak round2\_\allowbreak gate\_\allowbreak d8083dbc.txt}).}
\label{tab:floor}
\begin{tabular}{lrr}
\toprule
Feature & Re-run, scored rows & Gate test rows, clustered \\
\midrule
\texttt{velocity\_ratio\_1h\_vs\_30d} & \nReVelocity{} $\pm$\nReVelocityci & \nGrVelocity{} \nGrVelocityci \\
\texttt{counterparty\_is\_new\_for\_account} & \nReCounterparty{} $\pm$\nReCounterpartyci & \nGrCounterparty{} \nGrCounterpartyci \\
\texttt{tx\_count\_1h} & \nReTxHour{} $\pm$\nReTxHourci & \nGrTxHour{} \nGrTxHourci \\
\texttt{implied\_speed\_kmh} & \nReSpeed{} $\pm$\nReSpeedci & -- \\
\texttt{seconds\_since\_last\_tx} & \nReSeconds{} $\pm$\nReSecondsci & -- \\
\bottomrule
\end{tabular}
\end{table}

The velocity ratio, \texttt{tx\_count\_1h}, \texttt{implied\_speed\_kmh} and
\texttt{seconds\_since\_last\_tx} are burst or recency indicators; counterparty novelty is not. The
generator stages several scenarios as incidents of many rows: a SIM-swap drain, a velocity run and a
mule fan-out are each a rapid sequence of transactions on one account. We did not suppress that
structure, and we did not restate the control as a claim about columns only, because that would
redefine a control so that it passes. Whether real drains and fan-outs are as bursty as the
generator's is not established here. The consequence is a rule applied to every gate ranking metric in this paper, that is, to
the gate ranking metrics (AUC, recall at the false-positive budget and the per-channel AUCs): \textbf{such a metric is reported as a margin over the best single feature and the
best trivial rule, measured on the same rows}.

The floor was itself corrected once. The first measurement was taken before a review found that
\nUnboundedFixed{} features claiming unbounded history derived it from the truncated corpus they
were given. The correction moved \texttt{counterparty\_is\_new\_for\_account} from
\nCounterpartyNewBefore{} to \nCounterpartyNewM{} and left the velocity ratio unchanged; the review
had expected the headline to be inflated, and the measurement says it was not.
